% -------------------------------------------------------
%  پروژه پایانی سیستم‌های نهفته
%  سامانه هوشمند نگهبانی -- ماهان مجلسی (۴۰۲۱۷۰۵۱۶)
% -------------------------------------------------------

%% ============================================================
\section{معماری کلی سامانه}
%% ============================================================

سامانه از سه بخش تشکیل شده است که روی دو ماشین اجرا می‌شوند:

\begin{itemize}
	\item \textbf{برد}: یک \lr{Raspberry Pi 3 Model B+} با \lr{Ubuntu 24.04 LTS} (معماری \lr{aarch64}، ۹۰۰ مگابایت حافظه‌ی در دسترس). دیمن اصلی به زبان \lr{C} نوشته شده و وب‌سرور \lr{HTTPS}، \lr{REST API}، کلاینت \lr{MQTT}، ارسال ایمیل، جعبه‌ی سیاه \lr{SQLite}، \lr{watchdog} نرم‌افزاری و مدیریت حرارتی را بر عهده دارد.
	\item \textbf{ماژول پردازش تصویر}: یک فرایند مستقل \lr{Python} با \lr{OpenCV DNN} و مدل \lr{MobileNet-SSD} که تنها وظیفه‌اش تشخیص انسان است. طبق صورت پروژه، این تنها بخشی است که استفاده از \lr{Python} در آن مجاز است.
	\item \textbf{کامپیوتر شخصی}: \lr{Broker} پیام‌رسان \lr{Mosquitto} و همچنین منبع تصویر. چون برد دوربینی ندارد، طبق راهکار پیش‌بینی‌شده در صورت پروژه، تصویر وب‌کم لپ‌تاپ به برد استریم می‌شود.
\end{itemize}

\subsection{چرا دو فرایند جدا؟}

پردازش تصویر و منطق اصلی می‌توانستند در یک فرایند باشند، اما جدا نگه‌داشتن آن‌ها یک تصمیم آگاهانه است: مدل \lr{Caffe} حدود ۲۳ مگابایت است و بارگذاری آن چند صد میلی‌ثانیه طول می‌کشد، و مهم‌تر اینکه \lr{watchdog} بخش ۴--۳ باید بتواند \emph{فقط} آشکارساز را ری‌استارت کند بدون آنکه وب‌سرور، \lr{MQTT} و ایمیل قطع شوند. اگر هر دو در یک فرایند بودند، هر بازیابی دوربین کل سامانه را از دسترس خارج می‌کرد.

تبادل تصویر میان این دو از طریق یک ناحیه‌ی \textbf{حافظه‌ی مشترک} (\lr{/dev/shm/guardian\_frame}) انجام می‌شود. دیمن \lr{C} این ناحیه را می‌سازد و مالک آن است؛ فرایند \lr{Python} تنها به آن متصل می‌شود. انتخاب حافظه‌ی مشترک به‌جای سوکت یا لوله به این دلیل است که هر فریم \lr{JPEG} حدود ۱۵ تا ۳۵ کیلوبایت است و کپی‌کردن آن در هر تبادل، روی پردازنده‌ای که هر سیکل آن ارزشمند است، هزینه‌ی بی‌دلیلی است.

\subsection{مسیر تصویر}

\begin{latin}
\begin{verbatim}
  MacBook (webcam)                     Raspberry Pi 3B+
  ----------------                     ----------------
  ffmpeg avfoundation
    |  MJPEG over plain TCP :9000
    v
  guardian-ingest.socket  --->  guardian-vision (Python)
                                   |  detection + overlay
                                   v
                                /dev/shm/guardian_frame
                                   |
                                   v
                                guardian (C)
                                   |--> HTTPS :443  (dashboard, MJPEG, REST)
                                   |--> MQTT  ---> Mosquitto on the laptop
                                   |--> SMTP  ---> Gmail
                                   `--> SQLite  (black box)
\end{verbatim}
\end{latin}

انتقال تصویر عمداً ساده‌ترین شکل ممکن را دارد: تصاویر \lr{JPEG} پشت سر هم روی یک اتصال \lr{TCP} خام، بدون \lr{RTSP} و بدون \lr{container}. آشکارساز جریان ورودی را بر اساس نشانگرهای \lr{SOI/EOI} (\lr{FFD8}/\lr{FFD9}) به فریم‌ها می‌شکند. این کار تأخیر انتها-به-انتها را نزدیک به زمان یک فریم نگه می‌دارد که مستقیماً روی نتیجه‌ی آزمایش ۳--۵ اثر دارد.

نکته‌ی مهم این است که سوکت شنونده به \emph{یونیت جداگانه‌ای} به نام \lr{guardian-ingest.socket} تعلق دارد نه به خود آشکارساز. مزیت این کار آن است که پورت در لحظه‌ی رسیدن به \lr{sockets.target} باز می‌شود -- خیلی پیش از آنکه \lr{OpenCV} مدل را بارگذاری کند -- بنابراین لپ‌تاپ می‌تواند در هر زمانی متصل شود بدون آنکه با خطای \lr{connection refused} مواجه شود.

%% ============================================================
\section{ترتیب وابستگی سرویس‌ها}
%% ============================================================

صورت پروژه خواسته است که ترتیب وابستگی با \lr{After} و \lr{Requires} مدیریت شود و در نموداری با ذکر علت توضیح داده شود.

\begin{latin}
\begin{verbatim}
  sockets.target
       |
       v
  guardian-ingest.socket  (binds 0.0.0.0:9000, PartOf=guardian-vision)
       |  Requires= + After=
       v
  guardian-vision.service  (attaches to shm, retries up to 30 s)
       |  After=  (ordering only -- NOT Requires=)
       v
  guardian.service         (creates shm, HTTPS/API/MQTT/mail/blackbox)
       |  After=  + Requires=
       v
  guardian-api.service     (FastAPI/Swagger, proxies to 127.0.0.1:8081)
\end{verbatim}
\end{latin}

\subsection{یک اشتباه واقعی که در حین آزمایش‌ها کشف شد}

نسخه‌ی اولیه‌ی \lr{guardian.service} شامل \lr{Requires=guardian-ingest.socket} بود. استدلال ظاهراً درست بود: سوکت همان «نقطه‌ی آماده‌بودن دوربین» است. اما این وابستگی یک زنجیره‌ی خاموش می‌ساخت:

\begin{enumerate}
	\item سرویس \lr{guardian-vision} متوقف می‌شود.
	\item در پی آن \lr{guardian-ingest.socket} هم متوقف می‌شود، چون با \lr{PartOf} به آن گره خورده است.
	\item و در پی آن \lr{guardian.service} هم متوقف می‌شود، چون آن سوکت را \lr{Requires} کرده بود.
\end{enumerate}

یعنی متوقف‌کردن آشکارساز، وب‌سرور و \lr{MQTT} و ایمیل و جعبه‌ی سیاه را هم از بین می‌برد؛ و چون \lr{Requires} فقط توقف را منتشر می‌کند و چیزی دیمن را دوباره بالا نمی‌آورد، \textbf{هر بار که \lr{watchdog} بخش ۴--۳ آشکارساز را ری‌استارت می‌کرد، کل سامانه برای همیشه از کار می‌افتاد.} این ایراد در جریان آزمایش ۲--۱ مشاهده شد: دیمن با \lr{MainPID=0} غیرفعال بود و به همین دلیل تمام مقادیر حافظه‌ی ثبت‌شده صفر بودند.

راه‌حل این بود که این وابستگی به \lr{After=} تنزل داده شود. بررسی کد نشان داد که دیمن \lr{C} هیچ‌گاه به این سوکت دست نمی‌زند -- نه \lr{LISTEN\_FDS} دارد و نه هیچ ارجاعی به پورت ۹۰۰۰ -- پس این وابستگی از ابتدا واقعی نبود و تنها ترتیب اجرا لازم بود.

نکته‌ی ظریف دیگر آنکه \lr{guardian-vision} \emph{پیش از} \lr{guardian} اجرا می‌شود، در حالی که ناحیه‌ی حافظه‌ی مشترک را دیمن \lr{C} می‌سازد. این وارونگی عمدی و بی‌خطر است، چون آشکارساز تا ۳۰ ثانیه برای پیداشدن ناحیه تلاش می‌کند. بدین ترتیب نبود موقت یکی، دیگری را ساقط نمی‌کند.

%% ============================================================
\section{بخش اول: وب‌سرور، سرویس خودکار و \lr{SSL}}
%% ============================================================

\subsection{وب‌سرور}

وب‌سرور به زبان \lr{C} و بر پایه‌ی سوکت‌های \lr{POSIX} و \lr{OpenSSL} نوشته شده است. صفحه‌ی \lr{HTML} شامل عنوان با نام و شماره‌ی دانشجویی، استریم زنده، تعداد افراد تشخیص‌داده‌شده، و دما و حافظه و درصد مصرف پردازنده با به‌روزرسانی هر ۲ ثانیه است.

دما و حافظه \textbf{مستقیماً} از \lr{/sys/class/thermal/thermal\_zone0/temp} و \lr{/proc/meminfo} و \lr{/proc/stat} خوانده می‌شوند، نه با فراخوانی دستورات لینوکس. محاسبه‌ی درصد پردازنده از تفاضل دو نمونه‌ی متوالی \lr{/proc/stat} به‌دست می‌آید، چون این فایل شمارنده‌ی تجمعی از زمان بوت است و یک نمونه‌ی تنها معنایی ندارد.

\subsection{\lr{SSL} با گواهی \lr{Self-Signed}}

گواهی با \lr{openssl} ساخته شده و فیلد \lr{CN} آن شماره‌ی دانشجویی است. درخواست‌های \lr{http} با کد ۳۰۱ به \lr{https} هدایت می‌شوند.

انتخاب کد ۳۰۱ (دائمی) به‌جای ۳۰۲ (موقت) عمدی است: مرورگر پس از نخستین بازدید دیگر اصلاً درخواست رمزنگاری‌نشده ارسال نمی‌کند، پس پورت ۸۰ حداکثر یک‌بار برای هر کلاینت استفاده می‌شود.

گواهی علاوه بر \lr{CN}، فیلد \lr{SAN} هم دارد که آدرس \lr{IP} برد و نام‌های \lr{mahan} و \lr{mahan.local} و \lr{localhost} را پوشش می‌دهد. دلیلش این است که مرورگرهای امروزی برای تطبیق نام میزبان \emph{فقط} \lr{SAN} را می‌خوانند و \lr{CN} را نادیده می‌گیرند؛ پس گواهی به هر دو نیاز دارد: \lr{CN} برای برآورده‌کردن خواسته‌ی صورت پروژه و \lr{SAN} برای آنکه اصلاً قابل استفاده باشد.

\subsection{سرویس‌های \lr{systemd}}

هر سه برنامه با \lr{Restart=always} پیکربندی شده‌اند. دیمن اصلی از \lr{Type=notify} استفاده می‌کند و تنها زمانی \lr{sd\_notify(READY=1)} می‌فرستد که شنونده‌ی \lr{TLS}، پایگاه‌داده و نخ‌های کارگر همگی بالا آمده باشند. این باعث می‌شود \lr{After=guardian.service} در یونیت \lr{API} به معنای «پس از آنکه واقعاً بتواند پاسخ دهد» باشد، نه صرفاً «پس از \lr{fork}».

دیمن با کاربر غیرروت \lr{guardian} اجرا می‌شود و برای \lr{bind} روی پورت‌های ۸۰ و ۴۴۳ تنها از قابلیت \lr{CAP\_NET\_BIND\_SERVICE} استفاده می‌کند.

\subsection{نتایج آزمایش‌های بخش اول}

\subsubsection*{آزمایش ۱--۱: زمان بوت}

\begin{latin}
\begin{verbatim}
Startup finished in 11.812s (kernel) + 53.022s (userspace) = 1min 4.835s
graphical.target reached after 44.195s in userspace.

  2.617s  guardian.service
   115ms  guardian-ingest.socket
\end{verbatim}
\end{latin}

سهم سرویس‌های این پروژه از کل بوت ناچیز است. سنگین‌ترین یونیت‌ها متعلق به خود \lr{Ubuntu} هستند:

\begin{center}
\begin{tabular}{lr}
\hline
یونیت & زمان \\
\hline
\lr{snapd.seeded.service} & ۱۷٫۵۰۴ ثانیه \\
\lr{snapd.service} & ۱۵٫۸۲۵ ثانیه \\
\lr{cloud-config.service} & ۱۴٫۷۴۵ ثانیه \\
\lr{cloud-final.service} & ۸٫۷۵۶ ثانیه \\
\hline
\end{tabular}
\end{center}

نکته‌ای که باید در تفسیر این جدول رعایت شود: \lr{systemd-analyze blame} زمان \emph{مقداردهی هر یونیت} را نشان می‌دهد نه مسیر بحرانی را؛ این زمان‌ها هم‌پوشانی زیادی دارند و جمعشان برابر کل زمان بوت نیست. عددی که واقعاً توضیح می‌دهد سامانه از چه لحظه‌ای قابل استفاده شد، خروجی \lr{critical-chain} است:

\begin{latin}
\begin{verbatim}
guardian.service +2.617s
 \_ time-sync.target @31.360s
     \_ chrony.service @28.108s +3.182s
\end{verbatim}
\end{latin}

یعنی دیمن کند نیست، بلکه \emph{منتظر} است: حدود ۳۱ ثانیه صبر می‌کند تا ساعت با \lr{NTP} همگام شود. این وابستگی عمدی است، چون جعبه‌ی سیاه و پیام‌های \lr{MQTT} و ایمیل همگی برچسب زمانی دارند و بردی که پیش از تنظیم ساعت شروع به ثبت کند، رکوردهایی با تاریخ ۱۹۷۰ تولید می‌کند. حذف این وابستگی سرویس را حدود ۳۰ ثانیه زودتر بالا می‌آورد اما به قیمت از دست رفتن اعتبار برچسب‌های زمانی -- معامله‌ای که ارزشش را ندارد.

\subsubsection*{آزمایش ۱--۲: کشتن وب‌سرور با \lr{kill -9}}

\begin{center}
\begin{tabular}{lr}
\hline
\lr{PID} پیش از \lr{kill} & ۱۹۸۵۵ \\
\lr{PID} پس از ری‌استارت & ۲۱۷۱۱ \\
شمارنده‌ی ری‌استارت \lr{systemd} & افزایش یافت \\
زمان تا فرایند جدید & حدود ۱۰ ثانیه \\
\hline
\end{tabular}
\end{center}

سیگنال \lr{SIGKILL} قابل \lr{trap} کردن نیست، پس فرایند بدون هیچ پاک‌سازی‌ای از بین می‌رود. \lr{systemd} خرابی یونیت را می‌بیند و \lr{Restart=always} را با تأخیر \lr{RestartSec} اعمال می‌کند؛ به همین دلیل فاصله چند ثانیه است و نه آنی. افزایش شمارنده‌ی ری‌استارت همان چیزی است که ری‌استارت واقعی را از «اصلاً نمردن فرایند» متمایز می‌کند.

\subsubsection*{آزمایش‌های ۱--۴ تا ۱--۶}

\begin{center}
\begin{tabular}{lll}
\hline
شماره & آزمایش & نتیجه \\
\hline
۱--۴ & هدایت \lr{http} به \lr{https} & کد \lr{301} به \lr{https://192.168.100.26/} \\
۱--۵ & گواهی \lr{self-signed} & \lr{CN = 402170516}، \lr{subject == issuer} \\
۱--۶ & صفحه‌ی \lr{HTML} & شماره‌ی دانشجویی در عنوان و متن صفحه \\
\hline
\end{tabular}
\end{center}

\textbf{مشکلی که در حین آزمایش ۱--۵ پیش آمد:} دریافت گواهی با \lr{openssl s\_client} از روی \lr{macOS} همواره با خطای \lr{EHOSTUNREACH} شکست می‌خورد، در حالی که \lr{curl} به همان آدرس کد ۲۰۰ می‌گرفت. علت، سازوکار \lr{Local Network Privacy} در نسخه‌های اخیر \lr{macOS} است که دسترسی به شبکه‌ی محلی را به‌ازای هر فایل اجرایی مجوز می‌دهد، و باینری \lr{openssl} نصب‌شده با \lr{Homebrew} در فهرست مجاز نبود. راه‌حل: اجرای \lr{s\_client} روی خود برد. این کار علاوه بر رفع مشکل، از نظر روش‌شناسی هم بهتر است، چون گواهی را از سوکت زنده‌ی \lr{TLS} می‌خواند نه از فایل روی دیسک؛ بنابراین آنچه واقعاً \emph{ارائه} می‌شود آزموده می‌شود.

%% ============================================================
\section{بخش دوم: \lr{RESTful API} و مانیتورینگ زنده}
%% ============================================================

\subsection{طراحی \lr{API}}

تمام منطق در کد \lr{C} پیاده شده و لایه‌ی \lr{FastAPI} تنها نقش مستندساز و درگاه را دارد؛ این لایه درخواست‌ها را به \lr{127.0.0.1:8081} روی خود برد \lr{proxy} می‌کند و هیچ محاسبه‌ای انجام نمی‌دهد.

\begin{center}
\begin{tabular}{lll}
\hline
\lr{ENDPOINT} & \lr{METHOD} & عملکرد \\
\hline
\lr{/api/v1/stream} & \lr{GET} & استریم زنده‌ی تصویر (\lr{MJPEG}) \\
\lr{/api/v1/persons} & \lr{GET} & تعداد افراد فعلی + \lr{timestamp} \\
\lr{/api/v1/telemetry} & \lr{GET} & دمای \lr{CPU}، حافظه‌ی آزاد، درصد بار \\
\lr{/api/v1/command} & \lr{POST} & دستور اجرایی، مثلاً \lr{\{"cmd":"reboot"\}} \\
\lr{/api/v1/history} & \lr{GET} & تاریخچه‌ی تشخیص‌ها \\
\hline
\end{tabular}
\end{center}

\lr{endpoint} دستور اجرایی به‌صورت جدولی از دستورات پیاده شده است تا افزودن دستور تازه در آینده تنها یک سطر باشد. دستورات فعلی: \lr{ping}، \lr{guard\_on}، \lr{guard\_off}، \lr{set\_fps}، \lr{set\_resolution}، \lr{set\_net\_input}، \lr{snapshot}، \lr{publish\_telemetry}، \lr{restart\_vision} و \lr{reboot}. دستورهایی که وضعیت سامانه را تغییر می‌دهند نیازمند \lr{Bearer token} هستند و دستورهای خواندنی خیر.

\subsection{نتایج آزمایش‌های بخش دوم}

\subsubsection*{آزمایش ۲--۱: دمای پردازنده در سه حالت}

نمونه‌برداری هر ۳۰ ثانیه به مدت ۵ دقیقه در هر حالت:

\begin{center}
\begin{tabular}{lrrrrr}
\hline
حالت & دمای کمینه & بیشینه & میانگین & میانگین \lr{CPU} & میانگین \lr{FPS} \\
\hline
بیکار & ۵۶٫۹۲ & ۵۹٫۰۷ & ۵۷٫۷۰ & ۴٫۰۰٪ & ۰٫۰۰ \\
فقط استریم & ۵۶٫۳۸ & ۵۸٫۰۰ & ۵۷٫۱۶ & ۱۰٫۶۸٪ & ۴٫۸۰ \\
استریم + تشخیص & ۶۰٫۱۵ & ۶۳٫۳۸ & ۶۱٫۵۸ & ۴۳٫۱۷٪ & ۷٫۴۹ \\
\hline
\end{tabular}
\end{center}

سه نکته از این جدول بیرون می‌آید:

\textbf{یک.} انتقال تصویر عملاً رایگان است. تفاوت دمای «بیکار» و «فقط استریم» کمتر از یک درجه و در محدوده‌ی نوفه است، در حالی که مصرف پردازنده از ۴ به حدود ۱۱ درصد می‌رسد. یعنی رمزگشایی و رمزگذاری \lr{JPEG} و انتشار فریم بار محسوسی نیست؛ آنچه برد را گرم می‌کند شبکه‌ی عصبی است.

\textbf{دو.} حالت «فقط استریم» و «استریم + تشخیص» تقریباً \emph{نرخ فریم یکسانی} دارند (هر دو حدود ۷ فریم بر ثانیه در نمونه‌های پایدار). این قوی‌ترین شاهد ممکن بر درستی جداسازی استنتاج از مسیر فریم است: روشن یا خاموش بودن تشخیص دیگر بر روانی تصویر اثری ندارد. پیش از اصلاح معماری، همین دو حالت اختلاف چندبرابری داشتند.

\textbf{سه.} مصرف پردازنده در حالت تشخیص نسبت به پیش از اصلاح افزایش یافته است (۴۳ درصد در برابر ۲۹ درصد). این افزایش \emph{مطلوب} است نه نگران‌کننده: چون حلقه‌ی فریم دیگر منتظر استنتاج نمی‌ماند، نخ کارگر در واحد زمان فریم‌های بیشتری می‌گیرد و هسته‌های بیکار مصرف می‌شوند. بهای آن حدود ۴ درجه دمای بیشتر است.

\textbf{یادداشتی درباره‌ی اعتبار ستون بیکار:} در اجرای نخست این آزمایش، مصرف پردازنده در حالت بیکار ۱٫۳۰ درصد اندازه‌گیری شد و در اجرای دوم ۴٫۰۰ درصد. تفاوت ناشی از خود فرایند نظارت است: در اجرای دوم چند بار برای بررسی وضعیت به برد \lr{SSH} زده شد و هر ورود \lr{SSH} روی این سخت‌افزار کار محسوسی است. عدد واقعی حالت بیکار نزدیک‌تر به ۱٫۳ درصد است و ۴ درصد شامل اثر ناظر می‌شود. این نکته عمداً ذکر شده چون همان جنس خطایی است که پیش‌تر باعث آلودگی حالت بیکار توسط \lr{watchdog} شده بود.

\subsubsection*{آزمایش ۲--۲: مصرف حافظه و بررسی نشتی}

این آزمایش دو بار اجرا شد: یک‌بار روی دیمنی که تازه راه‌اندازی شده بود و یک‌بار روی دیمنی که از پیش گرم بود. مقایسه‌ی این دو، خودِ روشِ سنجش نشتی را روشن می‌کند.

\begin{center}
\begin{tabular}{lrr}
\hline
 & دیمن تازه & دیمن گرم \\
\hline
\lr{RSS} در آغاز & ۲۲۱۷۶ & ۱۹۵۲۸ \\
\lr{RSS} در پایان & ۲۲۳۳۶ & ۱۹۵۲۴ \\
تغییر خالص (کیلوبایت) & ‎+۱۶۰‎ & ‎−۴‎ \\
شیب گرم‌شدن (۴۰٪ اول) & ۱٫۳۲۴۱ & ۰٫۱۱۳۳ \\
شیب حالت پایدار (۶۰٪ آخر) & \textbf{۰٫۱۶۶۲} & \textbf{−۰٫۰۷۸۹} \\
برازش ساده روی کل بازه & ۰٫۵۹۸۸ & ۰٫۰۰۰۱ \\
\hline
\end{tabular}
\end{center}

\textbf{نشتی حافظه وجود ندارد.} در اجرای دوم شیب حالت پایدار حتی \emph{منفی} است، یعنی \lr{RSS} به‌جای رشد، اندکی کاهش می‌یابد.

اما رسیدن به این نتیجه خود درس روش‌شناختی داشت. در اجرای نخست، برازش خط بر \emph{کل} بازه عدد ۰٫۵۹۸۸ را می‌دهد که از آستانه‌ی هشدار عبور می‌کند و در نگاه اول «نشتی» به نظر می‌رسد. مشکل در خود آزمون بود: \lr{RSS} یک دیمن تازه در دقیقه‌ی نخست بالا می‌رود تا بافرهای \lr{JPEG} و حافظه‌ی نهان \lr{TLS} و صفحات \lr{SQLite} به اندازه‌ی کاری خود برسند و سپس متوقف می‌شود؛ برازش یک خط بر گرم‌شدن و حالت پایدار \emph{با هم}، آن شیب اولیه را چنان گزارش می‌کند که گویی تا ابد ادامه دارد.

ستون دوم جدول همین را تأیید می‌کند: روی دیمن گرم، چون اصلاً بازه‌ی گرم‌شدنی وجود ندارد، هر سه شیب با هم می‌خوانند و برازش ساده هم گمراه‌کننده نیست. یعنی خطای اجرای نخست در \emph{داده} نبود، در \emph{نحوه‌ی برازش} بود.

معیار درست، شیب در حالت پایدار است. ملاک تشخیص نشتی واقعی این است که شیب در هر دو نیمه \emph{یکسان} بماند؛ در هر دو اجرا شیب مسطح می‌شود.

تفاوت ۲۲ و ۱۹ مگابایتی مقدار پایه‌ی \lr{RSS} میان دو اجرا هم ناشی از نشتی نیست: دیمن میان دو اجرا ری‌استارت شده و فرایند تازه هنوز به اندازه‌ی کاری کامل خود نرسیده بود. هر دو مقدار بسیار پایین‌تر از سقف ۲۲۰ مگابایتی تعیین‌شده در یونیت هستند.

\subsubsection*{آزمایش ۲--۳: ۵۰ درخواست هم‌زمان}

\begin{center}
\begin{tabular}{lr}
\hline
درخواست‌ها / پاسخ \lr{HTTP 200} & ۵۰ / ۵۰ \\
تأخیر پایه (بدون بار) & ۰٫۰۹۹ ثانیه \\
کمینه / میانگین / بیشینه زیر بار & ۰٫۲۷۸ / ۰٫۹۹۴ / ۱٫۵۴۱ ثانیه \\
صدک ۹۵ & ۱٫۴۶۷ ثانیه \\
نسبت به تأخیر پایه & ۱۰٫۱ برابر \\
دمای پردازنده & ۶۵٫۵۳ به ۶۵٫۵۳ درجه (بدون تغییر) \\
\hline
\end{tabular}
\end{center}

تأخیر حدود ده برابر می‌شود اما هیچ درخواستی رد یا گم نمی‌شود. دیمن هر اتصال را از یک استخر نخ کوچک سرویس می‌دهد، پس ۵۰ ورودی هم‌زمان \emph{صف} می‌کشند نه اینکه به ۵۰ نخ منشعب شوند؛ به همین دلیل دنباله‌ی توزیع کشیده می‌شود در حالی که تعداد خطاها صفر می‌ماند. این رفتار مطلوب است: انشعاب به ۵۰ نخ روی بردی با ۹۰۰ مگابایت حافظه به \lr{swap} و بدتر شدن وضع منجر می‌شد.

نکته‌ی دیگر آنکه هر درخواست واقعاً \lr{/proc} و \lr{/sys} را دوباره می‌خواند و از حافظه‌ی نهان پاسخ داده نمی‌شود، پس کار به‌ازای هر فراخوانی حقیقتاً تکرار می‌شود.

\subsubsection*{آزمایش ۳--۳: تغییر رزولوشن ورودی شبکه}

این آزمایش پس از اصلاح معماری (انتقال استنتاج به نخ جداگانه) تکرار شد و نتیجه‌ی آن به‌طرز آموزنده‌ای با نتیجه‌ی پیش از اصلاح تفاوت دارد.

\begin{center}
\begin{tabular}{lrrrrr}
\hline
ورودی شبکه & \lr{FPS} & زمان استنتاج & دما & نرخ تشخیص & اطمینان \\
\hline
۳۰۰ پیکسل & ۷٫۳۲ & ۱۳۷۴ میلی‌ثانیه & ۶۴٫۹۹ & ۱۰۰٪ & ۰٫۹۵ \\
۲۲۴ پیکسل & ۶٫۹۷ & ۸۱۲ میلی‌ثانیه & ۶۵٫۵۳ & ۹۷٪ & ۰٫۹۶ \\
۱۶۰ پیکسل & ۶٫۴۱ & ۴۵۲ میلی‌ثانیه & ۶۲٫۳۰ & ۱۰۰٪ & ۰٫۹۱ \\
\hline
\end{tabular}
\end{center}

پیش از اصلاح معماری، \lr{FPS} با کوچک‌شدن ورودی به‌طور یکنواخت زیاد می‌شد (۵٫۱۹ و ۶٫۱۳ و ۶٫۶۸) چون نرخ انتشار مستقیماً به زمان استنتاج گره خورده بود. اکنون \lr{FPS} تقریباً \emph{مسطح} است و حتی اندکی در جهت معکوس حرکت می‌کند.

این دقیقاً همان چیزی است که از جداسازی انتظار می‌رفت: نرخ استریم دیگر تابع زمان استنتاج نیست، بلکه به مسیر فریم (رمزگشایی \lr{JPEG}، تغییر اندازه، \lr{overlay}، رمزگذاری) و سقف \lr{target\_fps} محدود می‌شود. زمان استنتاج همچنان با مربع طول ضلع ورودی متناسب است -- نسبت‌های $224^2/300^2$ و $160^2/300^2$ به‌ترتیب حدود ۵۶ و ۲۸ درصد را پیش‌بینی می‌کنند و مقادیر اندازه‌گیری‌شده کمی بالاتر از این پیش‌بینی‌اند، که سهم هزینه‌های ثابت هر فریم است.

کاهش جزئی \lr{FPS} در رزولوشن‌های پایین‌تر توضیح محتملی دارد که با داده‌ها سازگار است: با کوتاه‌شدن زمان استنتاج، نخ کارگر در واحد زمان \emph{بیشتر} فریم می‌پذیرد (از حدود ۰٫۷۳ به ۲٫۲ استنتاج بر ثانیه)، و این بار افزوده با مسیر فریم بر سر پردازنده رقابت می‌کند. تفاوت حدود ۱۲ درصد است و بخشی از آن می‌تواند نوفه‌ی حرارتی هم باشد.

\textbf{نتیجه‌گیری درباره‌ی نقطه‌ی بهینه:} پس از جداسازی، کاهش رزولوشن دیگر روانی تصویر را بهتر نمی‌کند؛ تنها \emph{تازگی کادرها} را افزایش می‌دهد، چون کادرها زودتر به‌روز می‌شوند. بنابراین انتخاب ۳۰۰ پیکسل به‌عنوان مقدار پیش‌فرض منطقی است: بالاترین دقت را دارد، \lr{FPS} استریم را کم نمی‌کند، و تنها بهایش تأخیر بیشتر کادر روی سوژه‌ی متحرک است. اگر کاربرد نیازمند دنبال‌کردن دقیق سوژه‌ی سریع باشد، ۲۲۴ پیکسل مصالحه‌ی بهتری است.

\textbf{هشدار درباره‌ی ستون‌های دقت:} در این اجرا حضور فرد در کادر به‌صورت صریح اعلام نشده بود، بنابراین اعداد نرخ تشخیص و اطمینان تنها \emph{نشانه} هستند و نه دقت به‌معنای دقیق. بدون داده‌ی مرجع نمی‌توان تشخیص درست را از مثبت کاذب تفکیک کرد.
